\section{An actual triangle obstruction to linking extension}
\label{sec:136-linking-triangle-obstruction}

\paragraph{Archive, inherited input, and exact scope.}
The source problem is cote136, physical p.~121, footer120,
scan121, with uniform hypotheses on physical pp.~117 and119,
footers116 and118, scans117 and119.  We use the already checked
\textup{ARCH-136-SKELETON-OBSTRUCTION}, specifically its
\(d=2\) case and equations (136.SO1)--(136.SO6).
We do not reprove or strengthen its prescribed-completion
obstruction.  The new assertion realizes the exact intersection
in (136.LK16) by the actual reduced-representable degree-three
repair of that source.  It shows that the old degree-two
vertex compatibility alone does not imply preservation under
this prescribed repair.  It is not a counterexample to the
original integral common-source question: the two targets
below are both the standard theory itself.

Work over \(R=\mathbb Z\) or \(\mathbb Z_{(p)}\), where
\(p\) is prime, and put \(S=\Gamma_R(T)\) and
\(J=\mathbb Q[T]/S\).  Write \(B\) for the standard
admissible simplicial PD theory.  Fix any \(m\geq3\).

\paragraph{The inherited source and its normalization.}
The cited source is \(A=B\otimes_R D\), where the generators
\(V,U\) of \(D\) are the two copies of
\(\operatorname{DK}(R[2])\), in forms two and three,
both of weight \(m\), with
\[
                 dV=pU,\qquad dU=0.
                                                               \tag{136.LT1}
\]
Thus \(A_0=S\), \(A_1=B_1\), \(H^0(A_n)=S\),
and \(H^1(A_n)=H^2(A_n)=0\) in every weight.
The source is rationally vertically exact and its horizontal
rows are acyclic for \(w>mq\).  Its two maps
\(f_0,f_1:A\rightrightarrows B\) are the identity on
\(B\), send \(U\) to zero, and satisfy
\[
 f_0(V)=0,\qquad f_1(V)=\eta,
 \qquad \eta=x_0^{[m-2]}e_1e_2\in B_2^{m,2},
 \qquad \int_{\Delta[2]}\eta=T^{[m]}.
                                                               \tag{136.LT2}
\]
Every face of \(V,U,\eta\) is zero.  All these assertions,
including the unit divided-power period in (136.LT2), are the
specified inherited input.  In particular the old degree-two
vertex linking differences are zero, since \(H^2(A)=0\).
The weight \(m\) may exceed any prescribed fixed native
weight guard.

\paragraph{ARCH-136-LINKING-TRIANGLE-OBSTRUCTION.}
Adjoin a free relative strict differential PD pair to \(A\)
on \(\Delta[2]/\operatorname{sk}_0\Delta[2]\):
\[
       |c|=2,\quad |t|=1,\quad \mathrm{wt}(c)=\mathrm{wt}(t)=m,
       \qquad dc=U,\qquad dt=V-pc.
                                                               \tag{136.LT3}
\]
Let the resulting simplicial algebra be \(E\).  Both old maps
extend to strict maps \(g_0,g_1:E\rightrightarrows B\).
For every such pair of extensions, the normalized linking
difference \(\delta:N_1K\to J\) of the new two-classes
is nonzero on \(N_1K\cap\partial(N_2L)\) in arbitrarily
large weights.  Here
\[
 L=(S/p)\langle m\rangle\otimes_R
       R[\Delta[2]]/R[\operatorname{sk}_0\Delta[2]],
 \qquad K=\ker\{L\longrightarrow H^3(A),\ c\longmapsto[U]\}.
                                                               \tag{136.LT4}
\]
Consequently the chain extension required in (136.LK16) fails
for all these choices.  Moreover no full-boundary version of
this pair can be installed in the old algebra with extensions
to both of the fixed old maps.

\begin{proof}
The reduced representable is degreewise free on the nonconstant
simplex maps to \([2]\).  The differential in (136.LT3)
is therefore a genuine free relative PD construction; its square
is zero by (136.LT1), and it respects all simplicial operators.
Its vertex generators are identified with zero.  Thus \(E_0=S\).
An explicit pair of extensions is
\[
             g_0(c)=g_0(t)=0,
             \qquad g_1(c)=0,\quad g_1(t)=\kappa,
             \qquad d\kappa=\eta.
                                                               \tag{136.LT5}
\]
Such \(\kappa\in B_2^{m,1}\) exists by the vertical
exactness of \(B_2\), and positive forms have zero vertex
restrictions.  The free relative property gives the asserted
strict simplicial extensions, without requiring \(\kappa\)
to vanish on the edges.

Use \(\partial=d_0-d_1+d_2\) for the horizontal triangle
boundary and set
\[
                 \nu=\partial c\in E_1^{m,2},
                 \qquad b=-\partial t\in E_1^{m,1}.
                                                               \tag{136.LT6}
\]
Since the old \(U,V\) have zero faces,
\(d\nu=0\) and \(db=p\nu\).  In simplicial degree one
there are exactly three new edge pairs, with
\(dc_e=0\), \(dt_e=-pc_e\), over \(A_1=B_1\).
The low-cohomology calculation of (136.LK8)--(136.LK9), or
the same direct calculation with these zero old inputs, gives
\[
 H^0(E_1)=S,\quad H^1(E_1)=0,\qquad
 H^2(E_1)=(S/p)\langle m\rangle^{\oplus3}.
                                                               \tag{136.LT7}
\]
In particular \([\nu]\) is the alternating vector of
the three independent edge classes and has order exactly \(p\).
The same low calculation in all simplicial degrees identifies
\(H^2(E)\) with \(K\), since \(H^2(A)=0\).
The generators here have prescribed zero vertices, so the
relative coefficient module is exactly (136.LT4).

For completeness, the coefficient complex has
\[
 N_1L=(S/p)\langle m\rangle^{\oplus3},\quad
 N_2L=(S/p)\langle m\rangle,\quad
 \partial(1)=(1,-1,1),\quad N_iL=0\ (i\ne1,2).
                                                               \tag{136.LT8}
\]
The map \(N_2L\to N_2H^3(A)\) sends its generator to
\([U]\) and is an isomorphism.  Indeed the integral
\(S\)-linear contraction of \(B_n\) onto \(S\) identifies
\(H^3(A_n)\) with
\((S/p)\langle m\rangle\otimes_R
\operatorname{DK}(R[2])_n\).  In forms at most three the
only positive part of \(D_n\) is the free map \(V\xrightarrow{p}U\).
It follows that \(N_2K=0\) and \(N_1K=N_1L\).
Thus the class (136.LT6) maps to the nonzero element
\(\partial(1)\) in the precise intersection
\(N_1K\cap\partial(N_2L)\); it is not a boundary from
\(N_2K\).

Now allow \emph{arbitrary} extensions of the two fixed old maps.
Write \(c_j=g_j(c)\in B_2^{m,2}\) and
\(t_j=g_j(t)\in B_2^{m,1}\).  Their equations are
\[
              dt_0=-pc_0,
              \qquad dt_1=\eta-pc_1.
                                                               \tag{136.LT9}
\]
Every form-two restriction to \(B_1\) is zero.  Therefore
\(g_j(\nu)=0\); in the linking formula (136.LK10) its
target one-form primitives may both be chosen zero.
The target images of \(b\) are \(-\partial t_j\).
They are closed one-forms on \(\Delta[1]\), so their
zero-form primitives have endpoint differences equal to their
integral periods.  Identify
\(\mathcal V_1(J)=J^2/\operatorname{diag}J\) with \(J\)
by terminal vertex minus initial vertex.  Stokes' identity then
gives
\[
 \begin{aligned}
 \delta([\nu])
   &=-\frac{1}{p}
           \int_{\partial\Delta[2]}(t_1-t_0)\pmod S\\
   &=-\frac{1}{p}
           \int_{\Delta[2]}(\eta-p(c_1-c_0))\pmod S\\
   &=-\frac{T^{[m]}}p\pmod S\ne0.
 \end{aligned}                                                \tag{136.LT10}
\]
Every integral divided-power top form in \(B_2\) has an
integral period in \(S\), by the inherited basis integration
formula.  This is why the arbitrary \(c_1-c_0\) term
disappears modulo \(S\).  Thus (136.LT10) is independent
of every allowed choice of the new target images, not only
of corrections of \(t_j\) by closed forms.

For \(w\geq m\), multiply (136.LT6) by \(T^{[w-m]}\).
The new relation remains \(d(T^{[w-m]}b)=pT^{[w-m]}\nu\),
and coefficient linearity yields
\[
       \delta\bigl([T^{[w-m]}\nu]\bigr)
          =-\binom wm\frac{T^{[w]}}p\pmod S.
                                                               \tag{136.LT11}
\]
Take \(w=m+p^a\), with \(p^a>m\).  The coefficient
\(\binom{m+p^a}{m}\) is one modulo \(p\), by
\((1+Z)^{p^a}=1+Z^{p^a}\) in \(\mathbf F_p[Z]\).
These weights are arbitrarily large.  The element in (136.LT11)
still belongs to the intersection in (136.LK16), since it is
the boundary of the corresponding coefficient multiple of the
triangle generator in \(N_2L\).  This proves the claimed
uniform-tail failure of that exact chain extension.

Finally suppose instead that the complete boundary values
\(c^\partial,t^\partial\) were chosen in the old algebra
\(A(\partial\Delta[2])\) and that a full-boundary attachment
with (136.LT3) admitted extensions to both old maps.  All
boundary two-forms are zero because \(A_1=B_1\), so
\(c^\partial=0\).  The old maps agree on the entire
boundary, hence the two target one-forms \(t_0,t_1\)
would have identical restrictions there.  By (136.LT9) and
Stokes this would imply
\[
                    T^{[m]}
                      =p\int_{\Delta[2]}(c_1-c_0),
                                                               \tag{136.LT12}
\]
which is impossible in the free coefficient module \(S^m\).
This excludes every choice of old full-boundary data under
these retained old maps.  It does not exclude replacing the
old source or changing its prescribed maps.
\end{proof}

\paragraph{Exact consequence and remaining gate.}
The graded-injective coefficient argument in the full-boundary
linking theorem is valid at its stated scope.  Here its missing
boundary-installation hypothesis fails, and the reduced repair
displays the failure as a nonzero linking value on an actual
triangle boundary.  Thus the intersection condition cannot be
deduced merely from \(H^0=S\), \(H^1=H^2=0\), rational
vertical exactness, uniform horizontal acyclicity, and agreement
of the old target maps on the first skeleton.  A special source
construction could retain further compatible data; this result
does not rule it out.  The original integral common-source
and terminality claims remain open.

\paragraph{Verification status.}
The author completed a whole saved-proof reread, checking the
actual reduced PD presentation, every extension equation, the
normalized coefficient complex and kernel, the low-cohomology
identification, the orientation and Stokes signs, invariance
under arbitrary target choices, the unbounded coefficient tail,
and the full-boundary impossibility.  Reviewer15 and root have completed whole saved-proof checks
with PASS.  Root also reread the entire inherited SO source,
with actual completion recorded centrally at 04:51:23.222473 UTC.
Reviewer15 contributed the SO normalization and period calculation.  The source theorem and its integration
normalization are inherited checked inputs, not new assertions
of this module.
